% !TEX root = ../main.tex
% ============================================================
% Results chapter, revised after the Greenlight meeting with Kees.
% Figures are referenced as if they are stored in figures/ch7/.
% ============================================================

% Compile-safe figure inclusion for this draft chapter. If a plot has not yet
% been copied into figures/ch7, LaTeX prints a labelled placeholder instead of
% aborting with a graphics error.
\providecommand{\resultsfigure}[2][]{%
    \IfFileExists{#2}{%
        \includegraphics[#1]{#2}%
    }{%
        \fbox{%
            \begin{minipage}[c][0.32\textheight][c]{0.92\linewidth}
                \centering
                Missing figure file:\\[0.5em]
                \texttt{\detokenize{#2}}
            \end{minipage}%
        }%
    }%
}

\chapter{Results}
\label{chapter:results}

This chapter evaluates whether learned frequency transfer can be used in a
Helmholtz solver. The main conclusion is intentionally precise. A learned
transfer model can predict high-frequency wavefields with small relative field
error, but a Krylov method is controlled by residuals, not by field error
alone. A good-looking predicted solution may still be a poor algebraic initial
guess, because the Helmholtz operator can amplify small field errors into large
residual components.

The chapter is organised as a single chain of evidence. First,
Section~\ref{sec:results-1d} uses a one-dimensional Dirichlet problem as a
controlled setting where the residual-amplification mechanism can be diagnosed
exactly in an analytical eigenbasis. Second, Section~\ref{sec:results-1d-pml}
adds a one-dimensional PML. This removes the simple symmetric eigenbasis, but
it tests the same residual-facing diagnostics and isolates the role of the
absorbing layer. Third, Section~\ref{sec:results-2d} carries the same
questions to the two-dimensional FD/PML setting used elsewhere in the thesis.
Across all three settings, the same distinction is maintained: field error,
true residual, and CSL-preconditioned residual are related but not
interchangeable quantities.

All main solver-facing results in this chapter use CSL damping
$\beta=0.3$. Earlier experiments with $\beta=0.1$ showed the same qualitative
ordering, but they are not used as the main comparison here. This convention is
kept fixed to avoid mixing numerical effects from the learned start with
changes in the classical preconditioner.

Throughout this chapter, the learned model is used as a \emph{warm start}
unless stated otherwise. It chooses the initial guess $x_0$. It does not
modify the operator and it does not replace the CSL preconditioner. The CSL
preconditioner still defines the preconditioned Krylov process; the learned
model changes only the initial residual supplied to that process.

This convention determines the interpretation of every table in the chapter.
The true residual $\|b-Ax\|_2/\|b\|_2$ measures whether the original discrete
Helmholtz equation is being solved. The preconditioned residual
$\|M^{-1}(b-Ax)\|_2/\|M^{-1}b\|_2$ measures the residual after the CSL
preconditioner has acted on it. A learned warm start may improve one of these
metrics without improving the other.

% ============================================================
\section{One-dimensional Dirichlet results}
\label{sec:results-1d}
% ============================================================

The one-dimensional Dirichlet experiment isolates the mechanism in the cleanest
possible setting. The problem is the homogeneous Dirichlet Helmholtz equation
on a grid with $N=512$ unknowns and no PML. The frequency pair is
\[
    \omega_L = 16,
    \qquad
    \omega_H = 32.
\]
With the sign convention
\[
    A_\omega = -D_{xx} - \omega^2 I,
\]
the eigenpairs are known analytically:
\begin{equation}
    \lambda_k(\omega)
    =
    \frac{4}{h^2}
    \sin^2\!\left(\frac{\pi k}{2(N+1)}\right)
    -
    \omega^2,
    \qquad
    v_k(j)
    =
    \sqrt{\frac{2}{N+1}}
    \sin\!\left(\frac{j k \pi}{N+1}\right).
    \label{eq:results-dirichlet-eigenpairs}
\end{equation}
The eigenvectors are Euclidean-normalised. This matters because modal
coefficients of field errors and residuals can then be compared directly.

The learned model in this section is a one-dimensional U-Net trained to map
$u_{16}$ to $u_{32}$. Its input consists of the real and imaginary parts of the
low-frequency field and, in the final Dirichlet model, the real and imaginary
right-hand-side channels. The training objective is the full-grid complex
relative $\ell^2$ field loss. The best validation loss is approximately
$5.37\times 10^{-3}$, so the model is a strong field predictor in this
controlled setting.

% ------------------------------------------------------------
\subsection{Field accuracy is not residual accuracy}
\label{subsec:results-field-vs-residual}
% ------------------------------------------------------------

Let $u_\star$ be the exact high-frequency solution and $x_0$ an initial guess.
The initial field error and residual satisfy
\begin{equation}
    e_0 = u_\star - x_0,
    \qquad
    r_0 = b - A_H x_0 = A_H e_0.
    \label{eq:results-error-residual}
\end{equation}
Expanding in the Dirichlet eigenbasis gives
\begin{equation}
    e_0 = \sum_{k=1}^{N} c_k(e_0)v_k,
    \qquad
    r_0 = \sum_{k=1}^{N} \lambda_k c_k(e_0)v_k.
    \label{eq:results-modal-amplification}
\end{equation}
Thus
\[
    c_k(r_0) = \lambda_k c_k(e_0).
\]
Small field-error coefficients in modes with large $|\lambda_k|$ can therefore
become large residual coefficients. This is the central mechanism behind the
apparently contradictory behaviour of the learned warm start.

\begin{table}[htbp]
\centering
\caption{One-dimensional Dirichlet initial-state diagnostics for the
$16\to32$ pair with CSL damping $\beta=0.3$. The raw U-Net is an accurate field
predictor, but it is not a safe true-residual warm start. The residual gate
keeps only modal components that reduce the residual.}
\label{tab:results-1d-initial-state}
\begin{tabular}{lccc}
\toprule
Initial guess & Field error & $\|r_0\|_2/\|b\|_2$
& $\|M^{-1}r_0\|_2/\|M^{-1}b\|_2$ \\
\midrule
Cold start & $1.000$ & $1.000$ & $1.000$ \\
Raw U-Net & $0.071$ & $12.81$ & $0.173$ \\
Residual gate & $0.066$ & $0.708$ & $0.147$ \\
Two gated cycles & $0.060$ & $0.681$ & $0.129$ \\
\bottomrule
\end{tabular}
\end{table}

Table~\ref{tab:results-1d-initial-state} is the most compact version of the
first result. The raw U-Net reduces the field error from $1.000$ to $0.071$,
but its true residual is $12.81\|b\|_2$. This is not a plotting artefact and
not a failure of the linear solver. It is exactly the amplification described
by~\eqref{eq:results-modal-amplification}. The same raw residual looks much
less severe after CSL preconditioning, with preconditioned residual $0.173$.
This explains why the raw learned start can still be useful inside
CSL-FGMRES, even though it is poor in the original residual norm.

The residual gate gives the clearest diagnostic answer. It lowers the true
initial residual from $12.81$ to $0.708$, below the cold-start value, while
also slightly improving the field error. The learned model therefore contains
useful solver-facing information, but the information must be selected. Used
raw, it contains harmful modal components; gated by residual improvement, it
becomes algebraically safe.

\begin{figure}[htbp]
    \centering
    \resultsfigure[width=\linewidth]{figures/ch7/01_true_residual_vs_iteration.png}
    \caption{True residual versus FGMRES iteration for the 1D Dirichlet
    $16\to32$ problem with $\beta=0.3$. Source plot:
    \texttt{01\_true\_residual\_vs\_iteration.png}. The raw U-Net begins with
    a much larger true residual than the cold start, while the residual-gated
    starts begin below the cold residual. After several iterations the curves
    enter a similar CSL-controlled convergence regime.}
    \label{fig:results-1d-dirichlet-true-residual}
\end{figure}

\begin{figure}[htbp]
    \centering
    \resultsfigure[width=\linewidth]{figures/ch7/03_field_error_vs_iteration.png}
    \caption{Field error versus FGMRES iteration for the same 1D Dirichlet
    experiment. Source plot:
    \texttt{03\_field\_error\_vs\_iteration.png}.
    The learned starts are much closer to the exact solution at iteration
    zero, but field error alone does not determine the residual seen by the
    original Helmholtz equation.}
    \label{fig:results-1d-dirichlet-field-error}
\end{figure}

Figures~\ref{fig:results-1d-dirichlet-true-residual}
and~\ref{fig:results-1d-dirichlet-field-error} answer a specific concern from
the Greenlight discussion: the residual plots can show apparent gaps because
the methods start at very different residual levels. The gaps are real. They
show different initial residual states, not inconsistent plotting. After a
small number of Krylov iterations, the curves begin to collapse toward similar
slopes because the subsequent convergence is governed mainly by the
CSL-preconditioned operator. The learned start improves the initial state; it
does not replace the Krylov process.

% ------------------------------------------------------------
\subsection{Residual gating and the size of the solver benefit}
\label{subsec:results-gating-benefit}
% ------------------------------------------------------------

The residual gate accepts a learned modal component only if that component
reduces the corresponding modal residual relative to the current baseline. In
the Dirichlet setting this can be done exactly because the eigenvectors in
\eqref{eq:results-dirichlet-eigenpairs} are known. This is deliberately a
one-dimensional explanatory tool. It is not proposed as a production method
for the two-dimensional PML problem.

The iteration improvement from gating is modest. In the present
per-iteration experiment, the raw U-Net, residual gate, and two-gated-cycle
starts all converge in nearly the same number of CSL-FGMRES iterations once
the early residual state is corrected. This should not be hidden, because it is
an important scientific point. The gate proves that useful learned components
exist and can make the initial state residual-safe. It does not, by itself,
change the spectrum of the preconditioned Krylov operator. Therefore the
iteration-count gain is small, while the diagnostic value is large.

The correct interpretation is therefore:
\[
\text{learned field prediction}
\quad\not\Rightarrow\quad
\text{solver-compatible initial state}.
\]
The implication becomes valid only after a residual-aware selection step.

% ------------------------------------------------------------
\subsection{Left and right preconditioning}
\label{subsec:results-left-right}
% ------------------------------------------------------------

Krylov convergence can be monitored in different residual norms depending on
how the preconditioner is inserted. Left preconditioning solves
\[
    M^{-1}A x = M^{-1}b,
\]
and therefore directly minimises the residual of the transformed system,
$M^{-1}(b-Ax)$. Right preconditioning solves
\[
    A M^{-1} y = b,
    \qquad x = M^{-1}y,
\]
and therefore more directly targets the true residual of the original system.

For this thesis the true residual
\[
    \frac{\|b-Ax_k\|_2}{\|b\|_2}
\]
is the headline convergence metric, because it measures whether the iterate
solves the original Helmholtz equation. The preconditioned residual
\[
    \frac{\|M^{-1}(b-Ax_k)\|_2}{\|M^{-1}b\|_2}
\]
is reported as an explanatory diagnostic, because it describes what the CSL
preconditioner presents to the Krylov process.

\begin{table}[htbp]
\centering
\caption{Left versus right CSL preconditioning for the 1D Dirichlet
$16\to32$ problem. The initial values are identical because the initial guess
is the same. By iteration 10, right preconditioning gives smaller true
residuals, while left preconditioning gives smaller preconditioned residuals.}
\label{tab:results-left-right}
\begin{tabular}{llccc}
\toprule
Method & Side & Initial true residual & Iteration-10 true residual
& Iteration-10 preconditioned residual \\
\midrule
Cold & Left & $1.000$ & $1.06\times10^{-2}$ & $2.68\times10^{-3}$ \\
Cold & Right & $1.000$ & $2.52\times10^{-3}$ & $4.53\times10^{-3}$ \\
Raw U-Net & Left & $13.19$ & $4.96\times10^{-2}$ & $7.56\times10^{-4}$ \\
Raw U-Net & Right & $13.19$ & $6.73\times10^{-4}$ & $1.48\times10^{-3}$ \\
Residual gate & Left & $0.722$ & $4.88\times10^{-3}$ & $6.42\times10^{-4}$ \\
Residual gate & Right & $0.722$ & $6.10\times10^{-4}$ & $1.27\times10^{-3}$ \\
\bottomrule
\end{tabular}
\end{table}

\begin{figure}[htbp]
    \centering
    \resultsfigure[width=\linewidth]{figures/ch7/01_left_vs_right_true_residual.png}
    \caption{True residual for explicit left- and right-preconditioned GMRES
    on the 1D Dirichlet problem. Source plot:
    \texttt{01\_left\_vs\_right\_true\_residual.png}. Right preconditioning
    reduces the true residual more directly, which makes it easier to
    interpret as a solver-facing convergence curve.}
    \label{fig:results-left-right-true}
\end{figure}

\begin{figure}[htbp]
    \centering
    \resultsfigure[width=\linewidth]{figures/ch7/02_left_vs_right_precond_residual.png}
    \caption{Preconditioned residual for the same left- and
    right-preconditioned GMRES comparison. Source plot:
    \texttt{02\_left\_vs\_right\_precond\_residual.png}. Left
    preconditioning gives smaller preconditioned residuals because that is the
    residual norm minimised by the transformed system.}
    \label{fig:results-left-right-precond}
\end{figure}

Table~\ref{tab:results-left-right} and
Figures~\ref{fig:results-left-right-true}--\ref{fig:results-left-right-precond}
show why both residual norms must be named explicitly. Left preconditioning is
not wrong, and right preconditioning is not universally superior. They optimise
different residuals. For thesis presentation, the safest convention is to show
the true residual as the primary solver-facing metric and the preconditioned
residual as the explanation of the Krylov behaviour.

The Dirichlet study therefore supplies the analytical mechanism for the rest
of the chapter. A learned field can be accurate, and it can even look benign
after CSL preconditioning, while still having a large true residual. Residual
gating shows that this failure is not caused by the network being useless; it
is caused by useful and harmful modal components being mixed in the same field
prediction.

% ============================================================
\section{One-dimensional PML as an intermediate test}
\label{sec:results-1d-pml}
% ============================================================

The one-dimensional PML experiment is the bridge between the clean Dirichlet
analysis and the two-dimensional FD/PML problem. It keeps the problem small
enough to inspect carefully, but it removes the simple symmetric Dirichlet
eigenbasis. For that reason this section reports residual and error curves
only. It deliberately does not attempt a full eigenvalue/eigenvector analysis
of the PML operator. The purpose is narrower and more practical: test whether
the learned field remains compatible with the full discrete operator when the
absorbing layer is included.

The most important comparison is between three learned starts. The
\emph{green raw} start is an older PML model used directly. The
\emph{green zero} start zeros the PML part of that output. The
\emph{flux-full} start is trained on the full flux-form PML output and is used
without zeroing the absorbing layer.

\begin{table}[htbp]
\centering
\caption{One-dimensional PML initial-state diagnostics for the $16\to32$ pair
with $\beta=0.3$. The flux-full model improves all three initial metrics,
whereas zeroing the PML strongly damages the true residual.}
\label{tab:results-1d-pml-initial-state}
\begin{tabular}{lccc}
\toprule
Initial guess & Field error & $\|r_0\|_2/\|b\|_2$
& $\|M^{-1}r_0\|_2/\|M^{-1}b\|_2$ \\
\midrule
Cold start & $1.000$ & $1.000$ & $1.000$ \\
Green raw & $0.257$ & $1.924$ & $0.339$ \\
Green zero & $0.521$ & $37.14$ & $0.753$ \\
Flux-full & $0.022$ & $0.539$ & $0.034$ \\
\bottomrule
\end{tabular}
\end{table}

\begin{figure}[htbp]
    \centering
    \resultsfigure[width=\linewidth]{figures/ch7/combined_1d_dirichlet_pml_true_residual.png}
    \caption{Side-by-side true-residual FGMRES curves for the two
    one-dimensional $16\to32$ experiments with CSL damping $\beta=0.3$.
    Left: in the Dirichlet problem, the raw U-Net starts with a much larger
    true residual than the cold start, while residual gating makes the learned
    start algebraically safe. Right: in the PML problem, the flux-full model
    starts below the cold residual, while zeroing the PML output produces a
    much larger initial residual. Together these curves show that learned
    transfer is useful only when its output is compatible with the residual
    seen by the discrete Helmholtz operator.}
    \label{fig:results-1d-dirichlet-pml-combined}
\end{figure}

\begin{figure}[htbp]
    \centering
    \resultsfigure[width=\linewidth]{figures/ch7/03_pml_field_error_vs_iteration.png}
    \caption{Field error versus FGMRES iteration for the 1D PML experiment.
    Source plot:
    \texttt{03\_pml\_field\_error\_vs\_iteration.png}.
    The flux-full model starts much closer to the exact full-grid PML solution
    than the older PML starts.}
    \label{fig:results-1d-pml-field-error}
\end{figure}

Figure~\ref{fig:results-1d-dirichlet-pml-combined} puts the two
one-dimensional lessons on the same page. In the Dirichlet case, the issue is
spectral: the raw learned field is close in relative $\ell^2$ error, but some
modal components are amplified by $A_H$ into a large true residual. The
residual gate removes this failure mode and gives an initial residual below
the cold-start value. In the PML case, the issue is boundary compatibility:
the absorbing layer is part of the discrete operator, so a field that is good
in the physical region can still be a poor algebraic state if the PML values
are inconsistent.

This result answers the practical PML question. The absorbing layer is part of
the discrete operator. It is therefore not safe to treat it as an output region
that can simply be discarded. In Table~\ref{tab:results-1d-pml-initial-state},
zeroing the PML increases the true residual from $1.924$ to $37.14$ for the
older green model. By contrast, the flux-full model, trained on the full PML
field, improves the field error, true residual, and preconditioned residual at
the same time.

The one-dimensional evidence now has two parts. The Dirichlet experiment shows
why field accuracy alone is not enough: modal errors can be amplified into
large residuals. The PML experiment shows that compatibility with the
absorbing layer is also solver-facing: the PML values are part of the equation
being solved, not an optional boundary decoration. The two-dimensional
experiments test whether these two lessons remain visible in the full FD/PML
setting.

% ============================================================
\section{Two-dimensional FD/PML results}
\label{sec:results-2d}
% ============================================================

The two-dimensional experiments test whether the two one-dimensional lessons
survive in the realistic finite-difference/PML setting. The first lesson is
that field error is not the same as residual error. The second is that the PML
must be learned as part of the operator-compatible solution field. The grid is
$512\times512$ with a PML of width 112 on each side, leaving a
$288\times288$ physical interior. The tested frequency pairs are
\[
    16\to32,
    \qquad
    32\to64,
    \qquad
    64\to128.
\]
Each pair uses $9600$ full-grid FD/PML samples with complex Gaussian sources.
The main learned model is a full-grid flux-trained \texttt{TransferUNet} with
\texttt{base\_ch=32} and five levels. Unlike older interior-focused models,
the PML region is included in the training target.

The field prediction results show that the model learns nontrivial transfer
maps for all pairs. The best full-grid validation errors are approximately
$0.240$, $0.321$, and $0.339$ for $16\to32$, $32\to64$, and $64\to128$,
respectively; the corresponding interior errors are approximately $0.190$,
$0.252$, and $0.276$. These numbers establish that the network has learned a
meaningful solution-field map. They do not, by themselves, establish solver
acceleration. The next question is therefore not only how small the validation
loss becomes, but what physical structure is being learned by the optimisation
process.

\begin{figure}[htbp]
    \centering
    \resultsfigure[width=\linewidth]{figures/ch7/training_64_128.png}
    \caption{Training and validation curves for the full-grid FD/PML
    \texttt{TransferUNet} on the $64\to128$ pair. This is the hardest tested
    pair and shows the behaviour of the full-grid complex relative $\ell^2$
    objective.}
    \label{fig:results-2d-training-64-128}
\end{figure}

% ------------------------------------------------------------
\subsection{Training behaviour and learned field structure}
\label{subsec:results-2d-learned-structure}
% ------------------------------------------------------------

The training curve in Figure~\ref{fig:results-2d-training-64-128} is included
for more than bookkeeping. It shows that the hardest tested transfer pair is
still being optimised under a field-level objective: the supervised loss asks
for a high-frequency solution field, not for a residual-minimising correction.
This matters for interpreting the rest of the chapter. The machine-learning
pipeline succeeds in learning organised wavefield structure, but the loss does
not by itself enforce compatibility with the Krylov residual.

A useful mechanistic diagnostic is provided by an earlier analytic
Green's-function CNN experiment. This diagnostic is not the main quantitative
FD/PML benchmark used in Table~\ref{tab:results-2d-beta03-precond}; it uses an
older $N=600$ dilated CNN trained on analytic multi-source fields. Its value is
explanatory. It shows what kind of representation a field-trained
frequency-transfer network can discover before any solver-facing residual
test is applied.

\begin{figure}[htbp]
    \centering
    \resultsfigure[width=\linewidth]{figures/ch7/fig_voronoi_prediction.png}
    \caption{Mechanistic diagnostic for an older analytic Green's-function
    CNN on a six-source $16\to32$ example. The prediction preserves
    source-centred outgoing wave packets, while the normalised error is
    structured rather than uniform. The Voronoi overlay marks the geometric
    partition induced by the source locations. The diagnostic model is not the
    final full-grid FD/PML model; it is used here to interpret the kind of
    source-wise structure learned by field-level frequency transfer.}
    \label{fig:results-voronoi-prediction}
\end{figure}

The diagnostic in Figure~\ref{fig:results-voronoi-prediction} suggests a
source-wise interpretation of the learned map. A multi-source Helmholtz field
can be written schematically as
\[
    u_\omega(x) = \sum_i a_i \phi_\omega(x;x_i),
\]
where $\phi_\omega(\cdot;x_i)$ is a Green's-function-like outgoing atom
centred at source location $x_i$. The network was not given the individual
atoms, the source labels, a sparsity loss, or a Voronoi loss. Nevertheless,
the prediction is consistent with a windowed approximation of the form
\[
    \widehat{u}_{\omega_H}(x)
    \approx
    \sum_i W_i(x) a_i \phi_{\omega_H}(x;x_i),
\]
where $W_i$ behaves like a smooth spatial window associated with source $i$.
The natural geometric model for these windows is the Voronoi tessellation of
the source locations: inside a cell, one source dominates the local outgoing
wave; near cell boundaries, coherent interference between sources becomes
important and the error increases.

\begin{figure}[htbp]
    \centering
    \resultsfigure[width=0.92\linewidth]{figures/ch7/fig_atom_decomposition.png}
    \caption{Physical atom-decomposition diagnostic for the same analytic
    setting. The figure separates the coherent high-frequency target from a
    Voronoi-windowed atom approximation. The comparison clarifies the
    mechanism suggested by Figure~\ref{fig:results-voronoi-prediction}: local
    source-centred propagation can be learned well, while global coherent
    interference is the harder part of the field.}
    \label{fig:results-voronoi-atoms}
\end{figure}

This observation strengthens the interpretation of the validation curves. The
network is not merely fitting noise or producing a blurred high-frequency
field. It appears to infer source locations from the low-frequency wavefield
and re-emit local outgoing waves at the target frequency. At the same time,
the diagnostic explains why field accuracy is not the end of the story. A
Voronoi-windowed approximation can be physically meaningful and visually
structured while still omitting coherent interference terms that matter after
the high-frequency operator is applied.

The methodological point is central to the chapter. The optimisation did not
stop at the supervised loss value. Training behaviour, visual field
diagnostics, physical atom decompositions, true residuals, and preconditioned
residuals are used together to determine what has been learned and what
remains missing. This is why the remainder of the section evaluates the same
learned fields as warm starts in CSL-FGMRES rather than treating validation
RelL2 as the final success metric.

The solver-facing evaluation uses CSL-FGMRES with $\beta=0.3$, $10$
independently generated test right-hand sides per pair, and a fixed budget of
$40$ iterations. No method reaches the convergence tolerance within this
budget, so the iteration count is reported as $41$ for every method; this
value denotes failure to converge within the $40$-iteration budget, not an
additional iteration. The relevant comparison is therefore the residual after
the fixed iteration budget.

Two residual norms are reported. The true residual,
\[
    \frac{\|b-Ax\|_2}{\|b\|_2},
\]
measures whether the iterate solves the original Helmholtz equation. The
CSL-preconditioned residual,
\[
    \frac{\|M_{\mathrm{CSL}}^{-1}(b-Ax)\|_2}
         {\|M_{\mathrm{CSL}}^{-1}b\|_2},
\]
measures the residual in the coordinates seen by the preconditioned Krylov
process. This distinction is essential: a warm start can improve the final
true-residual trajectory without being uniformly better in the
preconditioned residual norm.

\begin{table}[htbp]
\centering
\caption{Two-dimensional FD/PML warm-start evaluation with CSL damping
$\beta=0.3$, $10$ test samples, and $40$ FGMRES iterations. All entries are
sample means. Since no method reaches the tolerance within the budget, the
final residuals are the main solver-facing comparison.}
\label{tab:results-2d-beta03-precond}
\begin{adjustbox}{max width=\textwidth}
\begin{tabular}{llcccccc}
\toprule
Pair & Method & Full error & PML ratio
& Initial true res. & Initial prec. res.
& Final true res. & Final prec. res. \\
\midrule
$16\to32$ & Cold
& $1.000$ & -- & $1.00$ & $1.00$
& $2.05{\times}10^{-3}$ & $1.56{\times}10^{-3}$ \\
& Depth5 zero
& $0.570$ & $0.000$ & $15.76$ & $0.233$
& $9.34{\times}10^{-4}$ & $4.32{\times}10^{-4}$ \\
& Flux-full raw
& $0.238$ & $0.409$ & $1.49$ & $0.116$
& $4.22{\times}10^{-4}$ & $2.30{\times}10^{-4}$ \\
& Flux-full zero
& $0.573$ & $0.000$ & $15.58$ & $0.236$
& $9.21{\times}10^{-4}$ & $4.62{\times}10^{-4}$ \\
\midrule
$32\to64$ & Cold
& $1.000$ & -- & $1.00$ & $1.00$
& $3.83{\times}10^{-1}$ & $3.08$ \\
& Depth5 zero
& $0.590$ & $0.000$ & $11.25$ & $0.273$
& $4.37{\times}10^{-1}$ & $0.764$ \\
& Flux-full raw
& $0.352$ & $0.365$ & $1.93$ & $0.712$
& $3.28{\times}10^{-1}$ & $1.09$ \\
& Flux-full zero
& $0.587$ & $0.000$ & $10.55$ & $0.265$
& $4.21{\times}10^{-1}$ & $0.749$ \\
\midrule
$64\to128$ & Cold
& $1.000$ & -- & $1.00$ & $1.00$
& $4.33{\times}10^{-1}$ & $10.92$ \\
& Depth5 zero
& $0.616$ & $0.000$ & $6.40$ & $0.895$
& $5.22{\times}10^{-1}$ & $6.97$ \\
& Flux-full raw
& $0.276$ & $0.322$ & $1.86$ & $18.41$
& $3.79{\times}10^{-1}$ & $28.51$ \\
& Flux-full zero
& $0.542$ & $0.000$ & $5.93$ & $0.685$
& $4.66{\times}10^{-1}$ & $5.31$ \\
\bottomrule
\end{tabular}
\end{adjustbox}
\end{table}

Table~\ref{tab:results-2d-beta03-precond} gives the central
two-dimensional result. The full-grid model used directly,
\texttt{flux\_full\_raw}, gives the lowest final true residual for all three
frequency pairs. The improvement is largest for $16\to32$, where the final
true residual decreases from $2.05\times10^{-3}$ for the cold start to
$4.22\times10^{-4}$. The harder pairs show smaller but still consistent
improvements:
\[
    3.83\times10^{-1} \to 3.28\times10^{-1}
    \quad\text{for }32\to64,
\]
and
\[
    4.33\times10^{-1} \to 3.79\times10^{-1}
    \quad\text{for }64\to128.
\]
Thus, in the true residual norm, the learned full-grid warm start improves the
finite-budget CSL-FGMRES trajectory for every tested pair.

The preconditioned residual gives a more nuanced picture. For $16\to32$,
\texttt{flux\_full\_raw} improves both the initial and final preconditioned
residuals. For $32\to64$, it improves the initial preconditioned residual
relative to the cold start, but \texttt{flux\_full\_zero} has the lowest final
preconditioned residual. For $64\to128$, the raw full-grid warm start has a
large preconditioned residual, even though it still gives the best final true
residual. This means that the learned warm start is not yet a uniformly
residual-compatible preconditioned state. It helps the original residual
trajectory, but it should not be interpreted as a learned preconditioner.
This is exactly the distinction introduced by the one-dimensional experiments:
the learned field contains useful solver-facing information, but the way that
information appears depends on which residual norm is inspected.

\begin{figure}[htbp]
    \centering
    \resultsfigure[width=\linewidth]{figures/ch7/gmres_clean_64_128.png}
    \caption{Clean true-residual CSL-FGMRES convergence curves for the
    $64\to128$ FD/PML pair at $\beta=0.3$. The plot omits the old raw depth5
    model to keep the scale readable. The full-grid flux-trained warm start
    gives the lowest final true residual after the fixed iteration budget,
    although the improvement is modest for this hardest pair.}
    \label{fig:results-2d-gmres-clean-64-128}
\end{figure}

\begin{figure}[htbp]
    \centering
    \resultsfigure[width=\linewidth]{figures/ch7/2d_final_true_residual_bars.png}
    \caption{Final true residual after $40$ CSL-FGMRES iterations for the
    2D FD/PML experiments at $\beta=0.3$. The full-grid flux-trained warm
    start has the lowest final true residual for all three frequency pairs.}
    \label{fig:results-2d-final-true-bars}
\end{figure}

\begin{figure}[htbp]
    \centering
    \resultsfigure[width=\linewidth]{figures/ch7/2d_initial_residual_comparison.png}
    \caption{Initial true residual and initial CSL-preconditioned residual for
    the cold start and the raw full-grid learned warm start. The learned start
    is strongly favourable in preconditioned coordinates for $16\to32$,
    partially favourable for $32\to64$, and unfavourable for $64\to128$.}
    \label{fig:results-2d-initial-residual-comparison}
\end{figure}



\subsection{Adaptive convergence-to-tolerance check}
\label{subsec:results-2d-adaptive-convergence}

The fixed-budget experiment above compares the finite-budget residual at a
common iteration count. It does not, however, answer the practical question of
how many iterations are needed to reach a prescribed tolerance. For that
reason I also ran adaptive stopping experiments with the same FD/PML operator,
the same CSL damping $\beta=0.3$, and the true-residual stopping rule
\[
    \frac{\|b-Ax_k\|_2}{\|b\|_2} \leq 10^{-6}.
\]
The $16\to32$ pair was evaluated on ten test right-hand sides. The
$32\to64$ pair was evaluated on three test right-hand sides, because it
requires approximately five times as many Krylov iterations. For the hardest
$64\to128$ pair, I ran a one-sample long convergence probe with a cap of
$10\,000$ iterations. This last run is included to establish the scale of the
iteration count rather than to estimate a sample mean.

\begin{table}[htbp]
\centering
\caption{Adaptive two-dimensional FD/PML convergence with CSL damping
$\beta=0.3$ and stopping tolerance $10^{-6}$ in true relative residual. The
learned starts are zero-PML warm starts: the physical-domain prediction is
used, while the PML strip is set to zero before FGMRES. The $64\to128$ row is
a one-sample long probe, so it should be read as a scale estimate rather than
a sample mean.}
\label{tab:results-2d-adaptive-all-pairs}
\begin{adjustbox}{max width=\textwidth}
\begin{tabular}{llcccc}
\toprule
Pair & Method & Samples & Converged & Mean it. & Initial true res. \\
\midrule
$16\to32$ & Cold & 10 & 10/10 & 69.4 & $1.00\times10^{0}$ \\
& Depth5 zero-PML & 10 & 10/10 & 66.9 & $1.58\times10^{1}$ \\
& Base32 zero-PML & 10 & 10/10 & 67.0 & $1.42\times10^{1}$ \\
& Base48 zero-PML & 10 & 10/10 & 67.1 & $1.43\times10^{1}$ \\
\midrule
$32\to64$ & Cold & 3 & 3/3 & 347.3 & $1.00\times10^{0}$ \\
& Depth5 zero-PML & 3 & 3/3 & 343.0 & $1.14\times10^{1}$ \\
& Base32 zero-PML & 3 & 3/3 & 343.7 & $9.92\times10^{0}$ \\
& Base48 zero-PML & 3 & 3/3 & 343.0 & $9.35\times10^{0}$ \\
\midrule
$64\to128$ & Cold & 1 & 1/1 & 2073.0 & $1.00\times10^{0}$ \\
& Depth5 zero-PML & 1 & 1/1 & 2062.0 & $6.10\times10^{0}$ \\
\bottomrule
\end{tabular}
\end{adjustbox}
\end{table}

\begin{figure}[htbp]
    \centering
    \resultsfigure[width=0.82\linewidth]{figures/ch7/2d_adaptive_summary/2d_adaptive_iterations_by_pair.png}
    \caption{Adaptive CSL-FGMRES iterations required to reach true relative
    residual $10^{-6}$ for the 2D FD/PML experiments. The vertical axis is
    logarithmic because the iteration count grows from about $70$ iterations
    for $16\to32$ to about $2000$ iterations for the $64\to128$ long probe.
    Learned zero-PML warm starts give small but consistent iteration savings.}
    \label{fig:results-2d-adaptive-iters-all-pairs}
\end{figure}

\begin{figure}[htbp]
    \centering
    \resultsfigure[width=0.72\linewidth]{figures/ch7/2d_adaptive_summary/2d_adaptive_iteration_savings.png}
    \caption{Iteration saving of the best learned zero-PML warm start relative
    to the cold start in the adaptive 2D FD/PML runs. The absolute saving grows
    with frequency, but the relative saving decreases because the underlying
    Krylov problem becomes much harder.}
    \label{fig:results-2d-adaptive-savings}
\end{figure}

\begin{figure}[htbp]
    \centering
    \resultsfigure[width=\linewidth]{figures/ch7/2d_adaptive_16_32/2d_16_32_true_residual_convergence_mean.png}
    \caption{Mean true-residual convergence curves for the adaptive
    $16\to32$ FD/PML experiment. Curves show the sample mean over ten test
    right-hand sides, with shaded interquartile bands. This plot shows the
    residual decay directly; the learned zero-PML starts cross the tolerance
    a few iterations before the cold start on average.}
    \label{fig:results-2d-adaptive-16-32-curves}
\end{figure}

Table~\ref{tab:results-2d-adaptive-all-pairs} and
Figures~\ref{fig:results-2d-adaptive-iters-all-pairs}--\ref{fig:results-2d-adaptive-16-32-curves}
show that the adaptive convergence result is positive but modest. For
$16\to32$, all methods converge on all ten test problems. The best learned
start is the depth-five zero-PML variant, with a mean of $66.9$ iterations
compared with $69.4$ for the cold start, a saving of $2.5$ FGMRES iterations.
For $32\to64$, the iteration count rises to roughly $347$ iterations for the
cold start and roughly $343$ iterations for the best learned starts, a saving
of about $4.3$ iterations on the three-sample probe. For $64\to128$, the
one-sample long run reaches the same tolerance only after about $2073$
iterations for the cold start and $2062$ iterations for the depth-five
zero-PML start.

The trend is therefore clear. Learned zero-PML warm starts consistently shave
a few FGMRES iterations from the adaptive convergence run, but the relative
benefit becomes smaller as the frequency increases: about $3.6\%$ for
$16\to32$, about $1.2\%$ for $32\to64$, and about $0.5\%$ in the single
$64\to128$ long probe. This is consistent with the fixed-budget evidence:
learned starts improve the initial state, but they do not change the
CSL-preconditioned Krylov operator. The high-frequency difficulty therefore
remains essentially a property of the classical preconditioned system.

The adaptive result should also be read together with the PML discussion
above. The zero-PML starts used here are not a statement that the PML is
unimportant. They are a controlled way to use the learned physical-domain
field while avoiding unsafe learned values in the absorbing layer. In the
fixed-budget table, the full-grid raw flux model can give the lowest final
true residual, whereas in the adaptive convergence-to-tolerance runs the
zero-PML variants give the smallest iteration counts among the starts tested.
The more robust conclusion is that learned starts contain useful
solver-facing information, while the best way to inject that information into
CSL-FGMRES, especially in the PML region, is still not fully solved.

The PML comparison is also important. For every frequency pair,
\texttt{flux\_full\_raw} has a lower final true residual than
\texttt{flux\_full\_zero}. Therefore, once the model is trained on the full
FD/PML field, keeping the learned PML output is better than zeroing it in the
primary solver-facing metric. This supports the conclusion from the
one-dimensional PML experiment: the absorbing layer is part of the discrete
operator, and solver-compatible learning must include it.

The older depth5 model illustrates the opposite behaviour. It has useful
interior information, but the raw PML output is harmful, especially for the
harder pairs. For example, at $64\to128$ the old raw depth5 start has final
true residual $15.86$, far worse than the cold-start value $0.433$. Zeroing
the old PML output avoids this catastrophic behaviour, but it does not produce
the best final true residual. The full-grid flux-trained model is therefore
not merely a better field predictor. In the true-residual metric it is more
compatible with the FD/PML solve than the PML-zeroed alternative, although the
preconditioned residual shows that this compatibility is not complete.

The two-dimensional conclusion is deliberately cautious. The neural model is a
warm start: it chooses the initial guess $x_0$. The CSL preconditioner still
preconditions the linear system, and FGMRES still performs the iterative solve.
The learned model should therefore not be described as a replacement for CSL or
as a learned preconditioner. Its current role is to provide a better initial
state for a classical preconditioned Krylov method. In the tested 2D FD/PML
setting, this improves the finite-budget true residual, and keeping the
learned PML output is better than zeroing it. At the same time, the
preconditioned-residual diagnostics show that residual compatibility is not
fully solved. This is the precise two-dimensional analogue of the 1D message:
learned frequency transfer is useful, but it must be made residual-aware
before it can become a robust solver component.

% ============================================================
\section{Chapter summary}
\label{sec:results-summary}
% ============================================================

The results support five conclusions.

First, field accuracy is not solver accuracy. In the 1D Dirichlet problem, the
raw U-Net reduces the field error substantially, but its true initial residual
can be much larger than the cold-start residual. The analytical Dirichlet
eigenbasis explains why:
\[
    c_k(r_0)=\lambda_k c_k(e_0).
\]
Small field errors in modes with large $|\lambda_k|$ can become large
residual components.

Second, residual-aware selection is necessary. The residual gate reduces the
Dirichlet true residual below the cold-start value. This proves that the
learned model contains useful solver-facing information, but also that the
information is not uniformly safe. The residual plots with apparent gaps are
therefore understood: they reflect different initial residual states, not a
plotting error.

Third, PML compatibility is a solver issue. In 1D PML, the flux-full model
improves field error, true residual, and preconditioned residual
simultaneously, while zeroing the PML can strongly worsen the residual. In 2D
FD/PML, the same principle appears in the true-residual metric: the full-grid
flux-trained model gives the best finite-budget true residual for all three
tested pairs, and keeping its learned PML output is better than zeroing it.

Fourth, the 2D learned fields have interpretable physical structure. The
Voronoi-windowed diagnostic suggests that a field-trained CNN can infer
source-centred outgoing atoms from a superposed low-frequency field, even
without source labels or a source-separation loss. This is scientifically
useful because it explains why the field predictions are nontrivial; it is
also a limitation because local source-wise structure does not guarantee that
the omitted interference terms are residual-safe.

Fifth, the learned model is useful as a warm start but is not yet a learned
preconditioner. In the 2D fixed-budget experiments, full-grid FD/PML training
improves the finite-budget true-residual trajectory, but the
CSL-preconditioned residual is not uniformly improved. The adaptive
convergence-to-tolerance runs give small iteration savings across the tested
frequency pairs: $69.4$ to $66.9$ mean iterations for $16\to32$, about
$347.3$ to $343.0$ iterations for the $32\to64$ probe, and $2073$ to $2062$
iterations in the single $64\to128$ long run. These numbers support the same
interpretation: the learned field improves the initial state, but it does not
change the underlying CSL-preconditioned Krylov process.

Overall, learned frequency transfer should be viewed as a source of
solver-relevant information, not as a drop-in replacement for Krylov methods.
In the current experiments, the learned model changes $x_0$ and therefore
changes the initial residual. The CSL preconditioner still defines the Krylov
search process. The natural next step is to train or gate learned components
to improve residuals and corrections directly, not only solution fields.
